% TikZ 图：Day 30 后段专题综合判断
\begin{tikzpicture}[scale=0.92,>=Latex]
  \tikzset{
    box/.style={draw,rounded corners=3pt,thick,fill=gray!6,
      minimum width=3.05cm,minimum height=2.35cm,
      text width=2.62cm,align=center,font=\small},
    title/.style={font=\bfseries},
    ray/.style={->,thick,blue},
    level/.style={thick},
    note/.style={font=\scriptsize,align=center}
  }

  \node[box] (opt) at (0,0) {};
  \node[title] at (0,0.82) {光学};
  \draw[thick] (-1.05,0.18) -- (1.05,0.18);
  \draw[ray] (-1.05,-0.45) -- (-0.20,0.18);
  \draw[ray] (-0.20,0.18) -- (0.82,0.66);
  \draw[dashed,blue] (-0.20,0.18) -- (0.82,-0.30);
  \node[note] at (0,-0.87) {先画光路\\再判反射、折射或成像};

  \node[box] (therm) at (4.0,0) {};
  \node[title] at (4.0,0.82) {热学};
  \draw[->,thick] (3.0,-0.55) -- (3.0,0.55) node[left] {$p$};
  \draw[->,thick] (3.0,-0.55) -- (4.95,-0.55) node[below] {$V$};
  \draw[very thick,orange] (3.18,-0.35) -- (4.55,0.38);
  \fill[orange!20] (3.18,-0.55) -- (3.18,-0.35) -- (4.55,0.38) -- (4.55,-0.55) -- cycle;
  \node[note] at (4.0,-0.87) {区分状态量和过程量\\面积对应气体做功};

  \node[box] (modern) at (8.0,0) {};
  \node[title] at (8.0,0.82) {近代物理};
  \draw[level] (7.15,0.35) -- (8.85,0.35) node[right] {$E_3$};
  \draw[level] (7.15,-0.05) -- (8.85,-0.05) node[right] {$E_2$};
  \draw[level] (7.15,-0.48) -- (8.85,-0.48) node[right] {$E_1$};
  \draw[->,red,thick] (7.58,0.30) -- (7.58,-0.42);
  \node[red,font=\scriptsize] at (7.15,-0.05) {$h\nu$};
  \node[note] at (8.0,-0.87) {先看能量量子化\\再用守恒关系};
\end{tikzpicture}
